%%%%%%%%%%%%%%%%%%%%%%%%%%%%%%%%%%%%%%%%%%%%%%%%%%%%%%%%%%%%%%%%%%%%%%%%%%%%%%%%%%%%%%%
% Copyright (c) 2011 QUALCOMM Incorporated. All rights reserved. 
% Qualcomm Confidential and Proprietary
%%%%%%%%%%%%%%%%%%%%%%%%%%%%%%%%%%%%%%%%%%%%%%%%%%%%%%%%%%%%%%%%%%%%%%%%%%%%%%%%%%%%%%%
%%%%%%%%%%%%%%%%%%%%%%%%%%%%%%%%%%%%%%%%%%%%%%%%%%%%%%%%%%%%%%%%%%%%%%%%%%%%%%%%%%%%%%%
% This custom LaTeX template file includes introductory information required for 
% Chapter 1 in the PDF document and must reside in the /dox/latex folder. This file is 
% pulled into the primary LaTeX document build file, refman.tex, during the document 
% build. To ensure that the PDF produced using this template file meets the Qualcomm 
% document formatting standards, except where indicated below, please do not change 
% this file. For information about the Living Docs Project, go to: go/livingdocs
%
% Author: Lori Boyters, Living Docs Technical Lead. 
% LaTeX templates are maintained in Perforce by the Living Docs team. As updates are 
% made, Doxygen/LaTeX users will be notified to update their templates. If you would 
% like to be notified when updates occur, send email to livingdocs.support.
% 
% Updated 05/28/2010 Rev B Version 8.
% Updated 08/06/2010 Rev B Version 9 for Doxygen 1.5.9
% Updated 09/22/2010 Rev C Version 1 for Doxygen 1.7.0
% Updated 01/28/2011 Rev C Version 2 (see Release Notes for details)
%%%%%%%%%%%%%%%%%%%%%%%%%%%%%%%%%%%%%%%%%%%%%%%%%%%%%%%%%%%%%%%%%%%%%%%%%%%%%%%%%%%%%%%

%%%%%%%%%%%%%%%%%%%%%%%%%%%%%%%%%%%%%%%%%%%%%%%%%%%%%%%%%%%%%%%%%%%%%%%%%%%%%%%%%%%%%%%%%%%%
\chapter{Introduction} 
\rm

%%%%%%%%%%%%%%%%%%%%%%%%%%%%%%%%%%%%%%%%%%%%%%%%%%%%%%%%%%%%%%%%%%%%%%%%%%%%%%%%%%%%%%%
\section{Purpose}

This document describes the <Name of the API> (<Acronym of the API>) API. The <Name of the API> provides <include description here>.


%%%%%%%%%%%%%%%%%%%%%%%%%%%%%%%%%%%%%%%%%%%%%%%%%%%%%%%%%%%%%%%%%%%%%%%%%%%%%%%%%%%%%%%
\section{Scope}

This document is intended for software developers who will be using the <Acronym of the API> API.

This document provides the public interfaces necessary to use the features provided by the <Acronym of the API>) API. A functional overview and information on leveraging the interface functionality are also provided.


%%%%%%%%%%%%%%%%%%%%%%%%%%%%%%%%%%%%%%%%%%%%%%%%%%%%%%%%%%%%%%%%%%%%%%%%%%%%%%%%%%%%%%%
\section{Conventions}

Function declarations, function names, type declarations, and code samples appear in a different font. For example, \verb,#include,.

Code variables appear in angle brackets. For example, \verb,<number>,.

Commands and command variables appear in a different font. For example, \tt\textbf{copy a:*.* b:}\rm. 

Parameter directions are indicated as follows:

\begin{itemize} 
\item \space {\tt[in]} indicates an input parameter.   
\item \space {\tt[out]} indicates an output parameter.  
\item \space {\tt[in,out]} indicates a parameter used for both input and output. 
\end{itemize} 
\newpage 


%%%%%%%%%%%%%%%%%%%%%%%%%%%%%%%%%%%%%%%%%%%%%%%%%%%%%%%%%%%%%%%%%%%%%%%%%%%%%%%%%%%%%%%
\section{References}

Table \ref{tbl:refs} lists reference documents, which may include Qualcomm documents and non-Qualcomm standards and resources. Reference documents that are no longer applicable are deleted from this table; therefore, reference numbers might not be sequential.

\renewcommand{\thetable}{\thechapter-\arabic{table}} 
\setlength\LTleft\parindent %Left aligns table.
\setlength\LTright\fill     %Left aligns table.
\begin{longtable}[htb]{|p{.25in}|>{\raggedright}p{3.9in}|>{\raggedright}p{1.85in}|}
   \caption{Reference documents and standards} \label{tbl:refs} \vspace{-0.05in} \\
   \hline
   \multicolumn{1}{|c|}{\textbf{Ref.}} & \multicolumn{2}{c|}{\textbf{Document}} \\
   \hline
\endfirsthead
   \caption[]{Reference documents and standards (cont.)} \vspace{-0.05in} \\
   \hline
   \multicolumn{1}{|c|}{\textbf{Ref.}} & \multicolumn{2}{c|}{\textbf{Document}} \\
   \hline
\endhead
\endfoot
   \hline
\endlastfoot
\multicolumn{3}{|l|}{\textbf{Qualcomm}} \\
\hline
\label{Q1}Q1 & \textit{Application Note: Software Glossary for Customers} & CL93-V3077-1 \tn
\hline
\label{Q2}Q2 & \textit{Title} & 80-Xxxxx-x \tn
\hline
\multicolumn{3}{|l|}{\textbf{Standards}} \\
\hline
\label{S1}S1 & \textit{Title} & Standard Number (optional date) \tn
\hline
\label{S2}S2 & \textit{} &  \tn
\hline
\multicolumn{3}{|l|}{\textbf{Resources}} \\
\hline
\label{R1}R1 & \textit{Title} & ISBN No. \tn
%\hline (use only before additional rows; do not put \hline after last row)
\end{longtable} 


%%%%%%%%%%%%%%%%%%%%%%%%%%%%%%%%%%%%%%%%%%%%%%%%%%%%%%%%%%%%%%%%%%%%%%%%%%%%%%%%%%%%%%%
\section{Technical Assistance}

For assistance or clarification on information in this guide, submit a case to Qualcomm CDMA 
Technologies at \color{blue}{\underline{\url{https://support.cdmatech.com}}}\color{black}. 

If you do not have access to the CDMATech Support Services website, register for access or send 
email to \color{blue}{\underline{\href{mailto:support.cdmatech@qualcomm.com}{support.cdmatech@qualcomm.com}}\color{black}.  


%%%%%%%%%%%%%%%%%%%%%%%%%%%%%%%%%%%%%%%%%%%%%%%%%%%%%%%%%%%%%%%%%%%%%%%%%%%%%%%%%%%%%%%
\section{Acronyms}

For definitions of terms and abbreviations, see \hyperref[Q1]{[Q1]}. Table \ref{tbl:acronyms} lists terms that are specific to this document.

\renewcommand{\thetable}{\thechapter-\arabic{table}} 
\setlength\LTleft\parindent %Left aligns table.
\setlength\LTright\fill     %Left aligns table.
\begin{longtable}[Htb]{|>{\raggedright}p{2.5cm}|>{\raggedright}p{13.2cm}|}
   \caption{Acronyms} \label{tbl:acronyms} \vspace{-0.05in} \\
   \hline
   \multicolumn{1}{|c|}{\textbf{Acronym}} & \multicolumn{1}{c|}{\textbf{Definition}} \\
   \hline
\endfirsthead
   \caption[]{Acronyms (cont.)} \vspace{-0.05in} \\
   \hline
   \multicolumn{1}{|c|}{\textbf{Acronym}} & \multicolumn{1}{c|}{\textbf{Definition}} \\
   \hline
\endhead
\endfoot
   \hline
\endlastfoot
<AAA> & <description of AAA \tn
\hline
<BBB>& <description of BBB> \tn
%\hline (use only before additional rows; do not put \hline after last row)
\end{longtable}
